\section{Scalar errors, smoothing, and free enclosures}
\label{sec:scalar-error}

Fix the compact data and the scalar-flat attachment of
Proposition~\ref{prop:af-attachment}.  All solutions in this section are
the actual solutions, with homotopy parameter equal to one, of
Proposition~\ref{prop:compact-solutions}.  In particular, the notation
\(N_i,N_o,f,s,w,x,W,T,\ell,\Sigma_\pm,P\) has the meaning given in
Section~\ref{sec:compact-construction}.  Constants can depend on the fixed
compact data, the positive numbers \(b_*,\delta_0\), the chosen exterior
collar, and the fixed value of \(L\).  They are independent of \(M\), of
the sufficiently small subsequent lapse cutoff, and of the sufficiently
large finite penalty strengths, unless stated otherwise.  None of the
arguments in this section restricts the topology of the compact part.

We first state the output needed in the final comparison.  The capped
manifold in this statement is auxiliary; no extension of the original
initial data across its boundary is asserted.

\begin{proposition}[Smooth AF comparison and two enclosures]
\label{prop:af-comparison}
Let \((\Omega_i,g,K)\) and \((\Omega_o,g_o)\) be the fixed inner data
and attachment of Proposition~\ref{prop:af-attachment}, with attachment
mass \(m_o\).  Given \(0<\delta<1\) and positive tolerances
\(\eta_m,\eta_v,\eta_g\), the parameters of
Proposition~\ref{prop:compact-solutions} can be chosen so that there is
a smooth complete connected one-ended AF manifold \((X,h')\), a fixed
topological cap set \(O\) with \(\partial O=S\), and precompact smooth
open sets \(E_0,E\) containing \(O\) in their interiors, with the
following properties.
\begin{enumerate}[label=\textup{(\roman*)}]
\item The boundaries of \(E_0\) and \(E\) are compact embedded minimal
surfaces.  Both sets are outward minimizing, their open exteriors are
connected, and
\begin{equation}
 \Per_{h'}(E_0)\le \Per_{h'}(E).
 \label{eq:unsafe-two-perimeters}
\end{equation}
The set \(E_0\) is a largest perimeter minimizing enclosure of \(O\)
alone, is strictly outward minimizing, and every component of \(E_0\)
meets \(O\).
\item On the original side of the caps,
\begin{equation}
 \Scal_{h'}\ge-2,\qquad
 \Scal_{h'}\ge0\quad\hbox{on }X\setminus E,\qquad
 m_{h'}\le m_o+\eta_m .
 \label{eq:unsafe-comparison-output}
\end{equation}
In particular the closed exterior of \(E\) is smooth and complete,
has exactly one AF end, and has outer minimizing minimal boundary.
For any fixed \(1/2<\beta<1\), its end has metric decay
\(O_2(|x'|^{-\beta})\) and scalar decay
\(O(|x'|^{-2-2\beta})\).
\item Identify the two original base pieces with their images in
\(X\setminus\operatorname{int}O\), smoothly on each closed side of
the seam, and let \(g_{\rm base}\) denote their respective metrics
\(g,g_o\).  For the piecewise graph metric
\(h=x\bar g\), where \(\bar g=g+s^2\dd f^2\) inside and
\(\bar g=g_o\) outside, one has
\begin{equation}
 (1-\eta_g)h\le h'\le(1+\eta_g)h,\qquad
 \bar g\ge g_{\rm base},\qquad
 \dd V_{\bar g}=W\,\dd V_{\rm base},
 \label{eq:unsafe-metric-output}
\end{equation}
after decreasing \(\eta_g\) below one if necessary, and
\begin{equation}
 \int_{\{x<1-\delta\}}W\,\dd V_{\rm base}<\eta_v.
 \label{eq:unsafe-bad-volume-output}
\end{equation}
The scalar lower bound and the topology of \(X\setminus
\operatorname{int}O\) are independent of the large parameters.
\end{enumerate}
The order of choice is the following: choose \(b_*,\delta_0\) for the
mass tolerance, then \(\delta\), the thin exterior omitted collar and
\(L\) for the volume tolerance, then \(M\), then the lapse cutoff and
the finite penalty strengths.  The corner smoothing is chosen last,
separately for each resulting fixed solution.  In the enclosure
construction a large fixed scaled ball radius is chosen before \(M\).
\end{proposition}

We prove the proposition in several steps.  Throughout, a closed unsafe
set is used, so excluding it from the new exterior also excludes limits
of the regions in which the scalar sign is uncertain.

\subsection{Uniform geometry at unsafe points}

Define \(\mathcal U_M\subset\overline\Omega_i\) by the closed
versions of the four trace-cutoff inequalities in
\eqref{eq:compact-trace-cutoff}, together with \(N_i\ge\delta_0\).
The quotient at \(S\) has the continuous interpretation specified
there.  This is a compact set for each fixed solution.  The scalar
identity \eqref{eq:compact-scalar} and the definition of \(\Sigma_-\)
show that
\begin{equation}
 \Scal_h\ge0\quad\hbox{off }\mathcal U_M
 \quad\hbox{on each smooth side of the seam}.
 \label{eq:unsafe-scalar-support}
\end{equation}
Indeed the only possibly negative term is
\(-\ell^2(1-\zeta(N_i))\); it vanishes when the trace cutoff is
inactive or when \(N_i\le\delta_0\).

At \(P\in\mathcal U_M\), set
\begin{equation}
 N_0=N_i(P),\qquad r_P=x(P)^{-1/2}.
 \label{eq:unsafe-scale}
\end{equation}
Since \(f(P)\ge9M/10\), \(z(P)\le1/2\), and
\(w(P)\ge M^{9/10}\),
\begin{equation}
 x(P)\ge\frac{\sqrt3}{2}\delta_0 M^{9/10},\qquad
 r_P\le C M^{-9/20}.
 \label{eq:unsafe-radius-bound}
\end{equation}
The negative term in the scalar identity is bounded independently of
\(M\), so \eqref{eq:unsafe-scalar-support} also gives
\(\Scal_h\ge-1\) everywhere on the two smooth pieces for all
sufficiently large \(M\).

\begin{lemma}[Geometry at the unsafe scale]
\label{lem:unsafe-blowup}
There is a constant \(C\ge1\) with the following property.  For every
fixed \(R<\infty\), all sufficiently large \(M\), and every
\(P\in\mathcal U_M\), the base ball of radius \(Rr_P\) about \(P\)
does not meet the seam.  In interior charts, or in charts flattened at
\(S\) when necessary, whose unit of base length is \(r_P\),
\begin{equation}
 C^{-1}\Id\le h\le C\Id
 \label{eq:unsafe-uniform-metric}
\end{equation}
on the corresponding ball, restricted to the original side of \(S\).
The same comparison holds after reflection across \(S\).  The constant
\(C\) is independent of \(R\); the required lower bound for \(M\)
may depend on \(R\).  Equivalently, the order of the quantifiers is
\begin{equation}
 \exists C\ \forall R<\infty\ \exists M_0(R)\
 \forall M\ge M_0(R)\ \forall P\in\mathcal U_M.
 \label{eq:unsafe-quantifiers}
\end{equation}
\end{lemma}

\begin{proof}
We first control the normalized lapse without a slope bound.
Height compactness and a sublevel argument then exclude positive
limiting heights and a seam at bounded scaled distance.
Only after this do we use the gradient estimate to obtain the flat
graph limit. Finally a truncated Harnack argument controls the
conformal factor with a constant independent of the fixed scaled radius.

It is enough to examine an arbitrary sequence \(M\to\infty\) and
marked centers \(P\in\mathcal U_M\).  Write \(r=r_P\).  Use smooth
base charts, with Fermi charts at a boundary, and rescale their lengths
by \(r\); equivalently use the base metric \(g_r=r^{-2}g\).  On each
fixed scaled compact set the base metrics converge smoothly to the
Euclidean metric, or to the Euclidean metric on a half-space.  The two
compact boundary faces \(S\) and \(\mathcal C\) have positive base
separation, so at most one can remain at a bounded scaled distance.

\smallskip\noindent\emph{Lapse bounds before slope bounds.}
We first obtain lapse estimates without using height-gradient bounds.
Testing the height equation on nested base balls with a cutoff and a
height function bounded in absolute value by \(M\) gives, in scaled
volume units,
\begin{equation}
 \int_{B_H}(W^2-1)\,\dd V_{g_r}
 \le C_H\frac{M}{r}
       \int_{B_{H'}}N_i\,\dd V_{g_r},\qquad H<H'.
 \label{eq:unsafe-local-height-energy}
\end{equation}
For a ball reaching a Dirichlet face, subtract its constant boundary
height before testing; the boundary term then vanishes.  To verify the
estimate, use \(p\cdot\dd f=W^2-1\), \(|p|\le N_i\),
\(|T|\le C\), and a cutoff derivative of size \(C_H/r\).
These are also the only estimates required if the ball meets a face.

Put \(U=N_i/N_0\), and let \(A_{H'}\) be a supremum of this
normalized lapse on a larger fixed scaled patch.  The scalar forcing
in its rescaled equation obeys
\begin{align}
 \left\|\frac{r^2}{N_0}(\Sigma_++\Sigma_-)\right\|_{L^2(B_H,g_r)}
 &\le C_H\left(r^2+
       M^{1/2}r^{3/2}N_0^{-1/2}A_{H'}^{1/2}\right) \\
 &\le C_H\left(r^2+M^{-7/40}A_{H'}^{1/2}\right).
 \label{eq:unsafe-scaled-source}
\end{align}
Here \(\Sigma_+\le LW\), \(|\Sigma_-|\le C\delta_0\),
\(N_0\ge\delta_0\), and \eqref{eq:unsafe-local-height-energy} were
used.  The divergence drift has size \(O(r)\).  The natural flux
error at \(S\) has size \(O(rb_*/N_0)\); a bounded extension of the
normal datum incorporates it in the divergence forcing.

If no seam is present on the patch, the inhomogeneous Harnack estimate
for \(U\), including natural-flux reflection at \(S\), bounds the
smaller supremum by a constant times its infimum plus the right side
of \eqref{eq:unsafe-scaled-source}.  Choose the smaller patch to
contain the center, or use a boundary-foot ball containing it, so that
the infimum is at most \(U(P)=1\).  Thus, on finitely many nested
patches,
\begin{equation}
 A_H\le C_H\left(1+r^2+M^{-7/40}A_{H'}^{1/2}\right).
 \label{eq:unsafe-nonseam-iteration}
\end{equation}
The crude bound \(N_i\le C(L)M\) gives \(A_{H'}=O(M)\).
The exponent of \(M\) improves by
\(a\mapsto\max(0,a/2-7/40)\).  Two applications give
\(1\mapsto13/40\mapsto0\).  We have therefore obtained local
boundedness by a finite iteration, before asserting that the forcing
tends to zero.

There is a slightly different normalization argument if the seam has
bounded scaled distance.  Reflect the exterior lapse into the inner
side in the two base Fermi collars, and set
\(V=U+N_o^{\rm refl}/N_0\).  Use the inner density-cometric matrix
for the summed equation.  The common induced metric implies that the
two matrices differ by \(O(r|t|)\) at scaled distance \(|t|\) from
the face.  Positivity and the interior gradient estimate for the
harmonic exterior lapse on balls of radius comparable to \(|t|\)
give
\begin{equation}
 |(a_i-a_o^{\rm refl})\nabla(N_o^{\rm refl}/N_0)|
 \le Cr(N_o^{\rm refl}/N_0).
 \label{eq:unsafe-seam-drift}
\end{equation}
The exterior is harmonic on the fixed omitted collar, which contains
every such scaled patch for large \(M\).  The inner drift obeys the
same bound with \(U\).  Consequently the total correction is a
bounded divergence drift of size \(Cr\) times \(V\), with zero
natural flux at the face, and the scalar forcing is still bounded by
\eqref{eq:unsafe-scaled-source} with a larger supremum of \(V\).

The infimum in Harnack for \(V\) must be anchored to the inner
center value.  Here is that step explicitly.  Let \(m\) be the
infimum of \(V\) on a fixed smaller half-patch containing the
center and a central disk in the face.  On the face,
\(N_o=dN_i\), \(0<d\le1\), whence \(U\ge m/2\).
Choose a fixed smooth half-domain with that flat disk in its boundary
and containing the center.  The harmonic measure of a smaller disk
at the center is bounded below by a fixed positive number; this
remains true as the center approaches that disk.  Write \(U=H+e\),
where \(H\) is the harmonic replacement with the same boundary
values.  The zero-Dirichlet estimate for \(e\), using divergence
forcing in \(L^p\), \(3<p<6\), and scalar forcing in \(L^2\),
gives
\begin{equation}
 \|e\|_\infty\le
 C_H\left(rA_{H'}+r^2+M^{-7/40}A_{H'}^{1/2}\right).
 \label{eq:unsafe-replacement-error}
\end{equation}
The remaining boundary values of \(H\) are nonnegative, and
\(U(P)=1\).  Harmonic measure therefore bounds \(m\) by a
constant times one plus this error.  Combining this with the Harnack
estimate for the reflected sum yields
\begin{equation}
 A_H\le C_H\left(1+rA_{H'}+r^2+
                 M^{-7/40}A_{H'}^{1/2}\right).
 \label{eq:unsafe-seam-iteration}
\end{equation}
Starting again from \(O(M)\), the exponent improves by
\(a\mapsto\max(0,a-9/20,a/2-7/40)\).  Three nested applications
give \(1\mapsto11/20\mapsto1/10\mapsto0\).  Thus the sum is
locally bounded.  This argument uses no derivative of the seam trace
ratio and no separate trace compactness for either lapse.

The forcing in \eqref{eq:unsafe-scaled-source} now tends to zero.
Local Laplace estimates give \(W^{1,p}\) bounds, \(p>3\), and
H\"older compactness, including reflected natural boundary estimates
when appropriate.  Without a seam, every subsequential limit is a
positive harmonic function on all of \(\R^3\), or on a half-space
with zero Neumann data.  Positive harmonic Liouville, after reflection
in the latter case, and the center normalization show that the limit
is one.  At a seam the reflected sum has instead a positive constant
limit.  Its constant is positive because it is at least the inner
center value.  Its trace lower bound, the relation \(U\ge V/2\)
on the face, and harmonic replacement with the error
\eqref{eq:unsafe-replacement-error} tending to zero give a positive
lower bound for the inner lapse up to the face.  On compact open
inner sets it converges, in a smaller H\"older exponent, to a smooth
positive harmonic function.  These conclusions hold on every fixed
patch by diagonal extraction.  In particular both upper and positive
lower lapse bounds are now available in every case.

\smallskip\noindent\emph{Height compactness.}
We turn to the height, setting
\begin{equation}
 \mathcal F=\frac{N_0(M-f)}{r},\qquad
 \mathcal U=\frac{N_i}{N_0},\qquad \widetilde s=\frac{s}{N_0}.
 \label{eq:unsafe-normalized-height}
\end{equation}
It is nonnegative and \(\mathcal F(P)\le M^{-1}\) by the closed
cutoff inequality.  Its flux in the metric \(g_r\) is
\begin{equation}
 \mathfrak p=\widetilde s^2\nabla_{g_r}\mathcal F
       =-\frac r{N_0}p,\qquad
 \divg_{g_r}\mathfrak p=-rT\mathcal U=o(1).
 \label{eq:unsafe-height-equation}
\end{equation}
The constitutive rule is unchanged:
\(\widetilde s^2+|\mathfrak p|_{g_r}^2=\mathcal U^2\).
If \(\mathcal U\) is frozen at its actual value, the flux is
strictly monotone in its gradient argument \(\xi\), and the positive
upper and lower lapse bounds give
\begin{equation}
 \langle\mathfrak p(\xi),\xi\rangle
   \ge c\min(|\xi|^2,|\xi|),\qquad
 |\mathfrak p(\xi)|^2
   \le C\langle\mathfrak p(\xi),\xi\rangle.
 \label{eq:unsafe-flux-coercivity}
\end{equation}
For example, in an orthonormal frame
\(\mathfrak p(\xi)=a\xi\), where
\(a+a^2|\xi|^2=\mathcal U^2\); these assertions follow directly
by differentiation and by considering \(|\xi|\le1\) and
\(|\xi|\ge1\).

We require a compactness argument for possibly very steep graphs.
In the product of a scaled base chart with a height line, the bounded
vector field \((-\mathfrak p,1)\) has bounded divergence and its
pairing with the upward unit graph normal is bounded below by a
positive constant, by \eqref{eq:unsafe-flux-coercivity}.  Gauss--Green
applied to the subgraph and supergraph in product balls gives, for
either side's volume function \(V(a)\), an upper bound
\(C(V'(a)+V(a))\) for the graph perimeter in that ball.  It also
gives local perimeter bounds by using fixed surrounding balls.
The four-dimensional isoperimetric inequality, at sufficiently small
radii to absorb the volume term, yields
\(V'(a)\ge cV(a)^{3/4}\).  Thus at every interior graph point both
sides have volume at least \(ca^4\).  These constants are uniform
on each of the fixed patches under consideration.

At \(S\), \(\mathcal F=0\).  For product balls at heights bounded
away above zero, extend the subgraph as the empty set across the
face.  There is no new face perimeter, because its trace is empty
there.  The same subgraph density estimate is therefore valid even
when the ball reaches \(S\).  BV compactness on open inner product
patches and vertical nesting now give a subsequence such that all
bounded height truncations converge in local measure to those of an
extended measurable function \(\mathcal F_\infty\ge0\), whose
value is allowed initially to be \(+\infty\).

We next prove that a nonzero limiting height forces a positive lower
bound.  On a ball \(B_H\) contained in the inner side, test
\eqref{eq:unsafe-height-equation} with
\(\chi^2(k-\mathcal F)^+\), where \(\chi=1\) on \(B_h\).
By \eqref{eq:unsafe-flux-coercivity} and Young's inequality,
\begin{equation}
 \int_{B_h}\min(|\dd(k-\mathcal F)^+|^2,
                        |\dd(k-\mathcal F)^+|)
 \le C\left(\frac{k^2}{(H-h)^2}+\varepsilon_M k\right)
              |\{\mathcal F<k\}\cap B_H|,
 \label{eq:unsafe-sublevel-energy}
\end{equation}
where \(\varepsilon_M\to0\) on the fixed patch.  The same test is
allowed on a half-patch with extension by zero across a Dirichlet
face whose height exceeds \(k\).  On the portion with derivative at
most one, use Cauchy--Schwarz; on its complement use the linear part
of the energy.  The \(W^{1,1}\) Sobolev inequality for an
intermediate cutoff then gives, for \(0<l<k\),
\begin{equation}
 (k-l)|\{\mathcal F<l\}\cap B_h|^{2/3}
 \le C\left(\frac{k}{H-h}+
             \frac{k^2}{(H-h)^2}+o(1)\right)
                |\{\mathcal F<k\}\cap B_H|.
 \label{eq:unsafe-sublevel-sobolev}
\end{equation}
The error here includes \((\varepsilon_M k)^{1/2}\) and
\(\varepsilon_M k\), which tend to zero for fixed levels.

For completeness, take radii decreasing from \(H\) to \(H/2\)
and levels decreasing from \(k\) to \(k/2\) geometrically.  First
take \(M\) sufficiently large at the fixed \(H,k\) to absorb the
forcing into the first cutoff term.  This also absorbs it at every
later step, since the level remains at least \(k/2\) and the cutoff
derivatives only increase.  After dividing volumes by \(H^3\),
\eqref{eq:unsafe-sublevel-sobolev} has the form
\begin{equation}
 v_{j+1}^{2/3}\le C^{j+1}(1+k/H)v_j.
 \label{eq:unsafe-sublevel-iteration}
\end{equation}
For \(0<k\le H\), a sufficiently small initial fraction, with a
threshold independent of \(k/H\), therefore forces \(v_j\to0\).
If \(\mathcal F_\infty\) is not zero almost everywhere, choose
a density ball of one of its positive superlevel sets and then a
small fixed \(k>0\).  Convergence of truncations supplies that
small initial fraction.  Consequently \(\mathcal F\ge k/2\)
on a smaller open inner ball for all sufficiently large indices.

This lower bound propagates to compact subsets of the connected
limiting inner domain.  To justify propagation for the capped flux,
choose a fixed smooth connecting subdomain with a small nonnegative
Dirichlet bump on a boundary portion inside the known positive ball
and zero data elsewhere.  On this subdomain the actual lapse fields
converge in a H\"older exponent to smooth positive coefficients.
The flux is smooth and uniformly elliptic near zero gradient.  The
Dirichlet inverse of its linearization at zero maps divergence
\(C^\gamma\) data to \(C^{1,\gamma}\); contraction in a small
\(C^{1,\gamma}\) ball therefore supplies small-gradient barriers,
uniformly for the sequence, with divergence a positive constant
tending to zero that dominates the right side of
\eqref{eq:unsafe-height-equation}.  Their limiting solutions have
positive interior values by the strict maximum principle.  The
constructed boundary values are below those of \(\mathcal F\).
Strict gradient-flux monotonicity, testing the positive part of the
barrier minus \(\mathcal F\), proves comparison in the required
direction.  A finite chain of such domains and a finite covering
prove propagation on every compact inner subset.

In the \(S\) alternative, the lapse convergence is also H\"older
up to that face.  The connecting subdomains may then reach a flat
boundary patch, with zero data on that portion.  The limiting
barrier is smooth there and has strictly positive inward derivative
by Hopf's lemma.  The \(C^1\) convergence of the barriers gives
a uniform lower bound by a positive constant times inward scaled
distance.  This contradicts the closed unsafe condition
\begin{equation}
 \frac{\mathcal F(P)}{d_S(P)/r}\le M^{-1},
 \label{eq:unsafe-boundary-touchdown}
\end{equation}
with the normal derivative interpretation if \(P\in S\).
An interior limiting center is already excluded by
\(\mathcal F(P)\le M^{-1}\).  Thus a nonzero limiting height
is impossible in these two alternatives.

\smallskip\noindent\emph{Excluding seam limits.}
If the center approaches the seam, the boundary height is instead
\(N_0M/r\to\infty\).  A positive limiting height propagates
through the open half-space as above.  On a fixed ball reaching the
face, omit first a thin strip adjacent to the face.  The lower
bound on the remaining compact set, followed by a small choice of
\(k\), makes the starting sublevel fraction arbitrarily small.
Apply \eqref{eq:unsafe-sublevel-iteration} with the negative part
extended by zero across the seam.  It supplies a positive lower
bound up to the central face patch and contradicts center
touchdown.  This step does not require H\"older convergence of
the individual lapse traces.

The zero almost everywhere alternative is also impossible at a
seam.  On a fixed half-patch choose a smooth nonnegative function
\(b\) with amplitude \(k\), derivative \(O(k)\), positive by
\(ck\) on a central face disk and zero near the other edges.
The test \((b-\mathcal F)^+\) has zero boundary trace, since the
seam height tends to infinity.  Monotonicity and
\eqref{eq:unsafe-flux-coercivity} give
\begin{equation}
 \int_{\{\mathcal F<b\}}\min(|\dd\mathcal F|^2,
                                      |\dd\mathcal F|)
       \le Ck^2+o(1)k.
 \label{eq:unsafe-seam-energy}
\end{equation}
The function \(v=\min(\mathcal F,b)\) has trace \(b\) on
that disk and tends to zero in measure in the interior.  On almost
every line normal to a smaller central disk, the fundamental
theorem of calculus, followed by averaging over interior slices of
a fixed thin collar \(D_\varepsilon\), gives
\(\int_{D_\varepsilon}|\dd v|\ge ck-o(1)\).  Here convergence
in measure suffices because \(0\le v\le k\).  Since
\(\min(t^2,t)\ge t-1/4\) for \(t\ge0\), the left side of
\eqref{eq:unsafe-seam-energy} is at least
\begin{equation}
 ck-Ck|D_\varepsilon|-\tfrac14|D_\varepsilon|-o(1).
 \label{eq:unsafe-seam-trace-cost}
\end{equation}
Choose first \(k>0\) small enough that \(ck\) dominates
\(Ck^2\), and then the collar thin enough.  This contradicts
\eqref{eq:unsafe-seam-energy}.  All possible seam limits have now
been excluded.  Consequently
\begin{equation}
 \dist_g(P,\mathcal C)/r_P\longrightarrow\infty,
 \qquad \mathcal F_\infty=0
 \quad\hbox{almost everywhere in the remaining limits}.
 \label{eq:unsafe-seam-escape}
\end{equation}

\smallskip\noindent\emph{Uniform height and slope convergence.}
The height convergence is in fact uniform on every scaled compact
set, including up to \(S\).  Otherwise continuity and local
convergence in measure produce graph points at some fixed positive
height in a slightly larger patch.  A fixed small product ball at
such a point has uniformly positive subgraph volume by the density
bound proved above.  But that volume tends to zero because the
limiting height is zero.  If the points approach \(S\), use the
empty extension across the face; alternatively omit a thin boundary
strip before taking the limit.  The strip has too little product
volume to account for the density lower bound.  This excludes
positive spikes as well as bounded positive plateaus.

We can now apply the differential gradient estimate
Lemma~\ref{lem:gradient-maximum}.  Under the rescaling above the
synthetic endomorphism and trace target become \(-rK\) and
\(-rT\), and the lapse current becomes \((r^2/N_0)B\), so the
same identity applies.  Maximize
\begin{equation}
 -\log\widetilde s+j(\mathcal F)+k_0\log\chi,
 \label{eq:unsafe-local-gradient-test}
\end{equation}
where \(k_0\ge2\), \(\chi\) is a fixed local cutoff, and
\(j',j''\) are positive, bounded, and bounded away from zero on
the now bounded height range.  If \(m=|\dd\mathcal F|_{g_r}\),
the constitutive relations give
\begin{equation}
 qm^2=\frac{J}{\widetilde s^2}
       =\frac{JW^2}{\mathcal U^2},\qquad
 (1-y)^2(j'm)^2=\frac{(j')^2z^2}{\mathcal U^2}.
 \label{eq:unsafe-scaled-gradient-algebra}
\end{equation}
At an interior maximum with \(y\ge1/2\), the positive leading
term therefore controls \(W^2\).  The source term has size
\(r^2\Sigma_+/N_i\le CW\), and localization costs
\(C\chi^{-2}\).  Hence
\begin{equation}
 cW^2\le C(1+W+\chi^{-2}).
 \label{eq:unsafe-slope-bound}
\end{equation}
At \(S\), \eqref{eq:calculus-boundary-lapse}, \(B_n=-b_*\),
\(k_T\hat z\ge0\), and \(H_g\le0\) give
\(\partial_n(-\log\widetilde s)\ge0\) for the inward normal.
Also \(\partial_n\mathcal F=m\ge0\).  The derivative condition
at a boundary maximum in \eqref{eq:unsafe-local-gradient-test}
therefore gives \(j'm\le C/\chi\), which controls that case.
When \(y<1/2\), \(W\le\sqrt2\) directly.  The maximum test
and \(k_0\ge2\) thus bound \(W\) on a smaller patch.

After this finite slope bound, differentiating the height divergence
equation gives uniformly elliptic scalar divergence equations for
its first derivatives.  Their divergence forcing is controlled in
\(L^p\), \(p>3\), by the normalized lapse estimates.  Odd
reflection of \(\mathcal F\), with even reflection of lapse and
the Fermi metric, gives the same estimate at \(S\).  The resulting
gradient H\"older estimates and the uniform height convergence
imply
\begin{equation}
 \mathcal U\longrightarrow1,\qquad
 \dd\mathcal F\longrightarrow0,\qquad
 \widetilde s\longrightarrow1
 \label{eq:unsafe-flat-graph-limit}
\end{equation}
on every fixed scaled compact set.

\smallskip\noindent\emph{The conformal factor.}
It remains to control the conformal factor uniformly in the chosen
fixed radius.  Put \(w'=w/w(P)\).  On every fixed scaled patch,
the height just proved gives \(f\ge9M/10\), hence \(w\ge T_1\)
with \(T_1=M^{9/10}\).  The lapse is eventually at least
\(\delta_0/2\); choose the subsequent cutoff
\(\eps<\delta_0/8\).  Also
\(x=sw\to\infty\) there.  The volume source and the low-lapse
penalty are consequently absent.  The positive function \(w'\)
is a supersolution for its uniformly elliptic divergence principal
part, while
\(v'=(w'-a)^+\), \(a=4T_1/w(P)\), is a subsolution: above
that truncation the remaining height penalty vanishes as well.
Weak Harnack for \(w'\), combined with the local supremum bound
for \(v'\) with the same positive mean exponent, gives on balls
with fixed nesting ratios
\begin{equation}
 \sup_{B_R}w'
 \le a+C\left(\frac1{|B_{2R}|}\int_{B_{2R}}(w')^p\right)^{1/p}
 \le a+C'\inf_{B_R}w'\le(4+C')\inf_{B_R}w'.
 \label{eq:unsafe-w-harnack}
\end{equation}
The last step uses \(w'\ge T_1/w(P)=a/4\).  Natural zero flux
permits reflection at \(S\).

These constants can be chosen independently of the fixed radius.
Indeed, on each arbitrarily large fixed patch the normalized
principal matrices tend to the identity by
\eqref{eq:unsafe-flat-graph-limit}.  If the boundary is far enough,
use concentric balls about the center with universal radius ratios.
Otherwise use comparable larger balls about its boundary foot in a
single flattened chart, and reflect.  The smaller ball contains both
the center and the requested region.  Uniform ellipticity with fixed
constants and dilation invariance of the homogeneous estimates give
one common constant in \eqref{eq:unsafe-w-harnack}.  The index from
which that constant applies may depend on the radius.  Since
\(w'(P)=1\), it follows that \(C^{-1}\le w'\le C\).

In the rescaled coordinates the metric is
\begin{equation}
 h=\frac{x}{x(P)}
       \left(g_r+\widetilde s^2\dd\mathcal F^2\right),
 \qquad
 \frac{x}{x(P)}=\frac{\widetilde s}{\widetilde s(P)}w'.
 \label{eq:unsafe-rescaled-metric}
\end{equation}
Here, as usual, the expression denotes its pullback by the coordinate
rescaling.  Equations~\eqref{eq:unsafe-flat-graph-limit} and
\eqref{eq:unsafe-w-harnack} prove
\eqref{eq:unsafe-uniform-metric}.  At \(S\) the mixed metric
entries vanish, so reflection preserves this comparison.  If the
uniform statement of the lemma failed, a sequence violating it would
admit precisely the subsequences just analyzed, a contradiction.
\end{proof}

\subsection{Mass and a fixed enclosing comparator}

On the AF end outside a common compact support, \(N_o\) and
\(x=N_ow\) are harmonic and tend to one.  Radial barriers of
weight \(|x'|^{-\beta}\), \(1/2<\beta<1\), and estimates on
annuli rescaled to unit size give
\begin{equation}
 N_o=1+O_2(|x'|^{-\beta}),\qquad
 x=1+O_2(|x'|^{-\beta}).
 \label{eq:unsafe-af-decay}
\end{equation}
The constants in these end bounds can depend on the fixed solution.
Thus \(h=xg_o\) is AF of this order.  Since \(g_o\) is scalar
flat and \(x\) is harmonic there, its scalar curvature is
\(O(|x'|^{-2-2\beta})\), an integrable decay rate.

The ADM flux of the conformal metric error gives
\begin{equation}
 m_h-m_o=-\frac1{8\pi}
           \lim_{R\to\infty}\int_{|x'|=R}\partial_nx\,\dd A_{g_o}.
 \label{eq:unsafe-adm-flux}
\end{equation}
One may replace the conormal-area product by its Euclidean version:
the error is quadratic in the decaying fields and tends to zero.
On this end the current in \eqref{eq:calculus-compact-current} is
\(Q=N_o\nabla x+(1-x)\nabla N_o\), whose flux has the same
limit as \(\nabla x\).  Integration of its divergence over both
pieces, cancellation of the matched seam fluxes with their common
base area form, and the outward flux \(+b_*\) at \(S\) yield
\begin{equation}
 8\pi(m_h-m_o)
   =\int_S b_*\,\dd A_g
      +\int\big((x-1)\Sigma_++P-\Sigma_-\big)\,\dd V_{\rm base}.
 \label{eq:unsafe-mass-budget}
\end{equation}
This also proves existence of the ADM limit: the current divergence
has compact support.  By Lemma~\ref{lem:compact-penalties},
\begin{equation}
 \limsup_{M\to\infty}(m_h-m_o)\le C(b_*+\delta_0).
 \label{eq:unsafe-mass-limsup}
\end{equation}
The constant on the right uses only the fixed compact data.  The
penalty cutoff makes \((x-1)\Sigma_+\le0\), and the positive
penalty costs tend to zero after the subsequent choices specified
in Section~\ref{sec:compact-construction}.

We will need a competitor whose area remains bounded as \(M\)
increases.  Fix a sphere \(\mathcal S_1\) in the attached end
beyond all source supports.  Let \(\psi_1\) be its exterior
capacitary potential: it is harmonic, is one on \(\mathcal S_1\),
and tends to zero at infinity.  Existence follows by solving on
compact annuli, using the preceding barriers, and taking a limit.
The function \(\gamma_1=-\partial_n\psi_1\) is strictly positive
on the sphere by Hopf's lemma.  Green's formula for \(x-1\) and
\(\psi_1\) gives
\begin{equation}
 \int_{\mathcal S_1}\gamma_1(x-1)\,\dd A_{g_o}
       =-\lim_{R\to\infty}\int_{|x'|=R}\partial_nx\,\dd A_{g_o}
       =8\pi(m_h-m_o).
 \label{eq:unsafe-capacity-comparator}
\end{equation}
The boundary term at infinity in the first Green identity vanishes
by \(\beta>1/2\); alternatively apply the harmonic flux identity
for \(x\) on the intervening annulus.  Positivity of \(x\), the
positive minimum of \(\gamma_1\), and the upper mass bound imply
\begin{equation}
 \Area_h(\mathcal S_1)
       =\int_{\mathcal S_1}x\,\dd A_{g_o}\le C_*.
 \label{eq:unsafe-comparator-bound}
\end{equation}
For the fixed earlier choices, \(C_*\) is independent of large
\(M\) and of the later parameters.  It is also independent of
the scaled radius of the obstacles introduced below.

\subsection{Removing the corner}

The collar mollification and conformal correction below use Miao's
corner-smoothing method \cite{Miao2002}. We give the local argument for
the present unsafe-set and cap construction, rather than invoke his
global positive mass theorem.

\begin{lemma}[Corner smoothing with controlled mass]
\label{lem:corner-smoothing}
For each sufficiently large fixed solution, the joined metric \(h\)
admits smooth AF replacements \(h'\) on the original side of the
caps such that, with errors made arbitrarily small separately at
that solution,
\begin{equation}
 h'=(1+o(1))h\ \hbox{as quadratic forms},\qquad
 m_{h'}=m_h+o(1),\qquad H_{h'}(S)<0.
 \label{eq:unsafe-smoothing-estimates}
\end{equation}
They satisfy \(\Scal_{h'}\ge0\) off \(\mathcal U_M\) and
\(\Scal_{h'}\ge-2\) everywhere on the original side of the
caps.  Their AF metric and scalar decay are those of
\eqref{eq:unsafe-af-decay}.
\end{lemma}

\begin{proof}
By Lemma~\ref{lem:unsafe-blowup} the closed unsafe set is disjoint
from the seam for large \(M\).  Glue the two metrics using their
own Gaussian normal collars.  In the resulting smooth manifold
structure the joined metric is continuous and has the form
\(\dd t^2+\lambda(t)\), piecewise smooth, with the inner side
at \(t<0\).  The common induced metric and the mean-curvature
inequality in \eqref{eq:compact-mean-curvature} give
\begin{equation}
 \frac12\tr_{\lambda(0)}
              \big(\lambda'(0^+)-\lambda'(0^-)\big)\le0.
 \label{eq:unsafe-corner-sign}
\end{equation}
This change of collar identification retains a smooth identification
of each original closed piece with its image; comparisons with the
base metric will always use those identifications.

Mollify \(\lambda\) in \(t\) with a nonnegative smooth even
kernel of radius \(a\), and interpolate back to the original
smooth metrics in \(2a\le|t|\le3a\).  The resulting \(h_a\)
converges uniformly to \(h\); its first normal derivatives and
fixed tangential derivatives are bounded.  Its second normal
derivative is the jump of \(\lambda'\) times the mollifier plus
a uniformly bounded remainder.  The interpolation has the same
bound, since on those smooth one-sided regions the mollification
error is \(O(a^2)\) before taking the two cutoff derivatives.
In the hypersurface scalar-curvature formula the only potentially
unbounded term is \(-\tr_\lambda\lambda''\).  Replacing the
inverse metric in its jump term by \(\lambda(0)^{-1}\) costs
\(O(a)O(a^{-1})=O(1)\); the leading term has favorable sign by
\eqref{eq:unsafe-corner-sign}.  Thus \(\Scal_{h_a}\ge-C(h)\)
in the smoothing collar, uniformly as \(a\downarrow0\).
Choose a smooth \(c_a\ge0\), bounded by a constant depending
on the fixed metric, supported in \(|t|<4a\), and majorizing
the negative scalar part created in that collar.

Each orientable component of \(S\) bounds a compact handlebody.
Attach such caps, extend the metric smoothly through \(S\), and
denote their union by \(O\).  This gives a smooth complete
one-ended AF manifold on which to solve the correction.  Its
homogeneous Sobolev inequality
\begin{equation}
 \|v\|_{L^6}\le C\|\nabla v\|_{L^2},
       \qquad v\in C_c^\infty(X),
 \label{eq:unsafe-af-sobolev}
\end{equation}
has a constant uniform for small \(a\).  Indeed all these metrics
are uniformly equivalent to one fixed smooth AF metric.  The
Euclidean Sobolev inequality on the end and compact local Sobolev
inequalities leave only a compact \(L^2\) term.  Integration
along radial rays, followed by Cauchy--Schwarz with the integrable
weight \(r^{-2}\), controls the trace at a fixed end sphere by
the end gradient energy.  Poincar\'e's inequality with this trace
controls the remaining compact term, proving
\eqref{eq:unsafe-af-sobolev}.

Solve in the completion for the gradient norm
\begin{equation}
 -8\Delta_{h_a}u_a=c_a(1+u_a).
 \label{eq:unsafe-conformal-correction}
\end{equation}
Since \(\|c_a\|_{3/2}\to0\), the bilinear form
\(8\int|\nabla v|^2-\int c_av^2\) is coercive by
\eqref{eq:unsafe-af-sobolev}.  The right side is a bounded
functional with norm tending to zero, because
\(\|c_a\|_{6/5}\to0\).  Lax--Milgram gives a solution with
\(\|u_a\|_6+\|\nabla u_a\|_2=o(1)\).  Testing with its
negative part and using coercivity proves \(u_a\ge0\).
Local boundedness for scalar divergence equations, treating the
bounded \(c_a u_a\) term as a zeroth-order coefficient and the
\(L^2\)-small \(c_a\) as forcing, gives
\(\|u_a\|_\infty=o(1)\) on a fixed compact set.  On the
unchanged end, \(u_a\) is harmonic and tends to zero: this
follows first from its \(L^6\) membership and scaled local
estimates.  Exterior comparison then gives the global supremum
bound.  Elliptic estimates imply smooth convergence to zero on
compact sets off the seam.  The same end barriers as before give
\(u_a=O_2(|x'|^{-\beta})\).

Set \(h'=(1+u_a)^4h_a\) and restrict to the original side of
the caps.  The conformal scalar formula is
\begin{equation}
 \Scal_{h'}=(1+u_a)^{-5}
       \big(-8\Delta_{h_a}(1+u_a)+\Scal_{h_a}(1+u_a)\big)
       =(1+u_a)^{-4}(\Scal_{h_a}+c_a).
 \label{eq:unsafe-corrected-scalar}
\end{equation}
It is nonnegative off the unsafe set.  On that set the metric
was not mollified, \(c_a=0\), and \(\Scal_h\ge-1\);
since \(u_a\ge0\), the lower bound remains at least \(-1\).
In particular the claimed fixed lower bound \(-2\) holds.
The end scalar decay is preserved by the conformal formula.

The smoothing itself changes no asymptotic metric.  The conformal
ADM flux, followed by integration of
\eqref{eq:unsafe-conformal-correction} on the capped manifold,
gives
\begin{equation}
 m_{h'}-m_h
   =-\frac1{2\pi}\lim_{R\to\infty}
                   \int_{|x'|=R}\partial_nu_a\,\dd A_{h_a}
   =\frac1{16\pi}\int_X c_a(1+u_a)\,\dd V_{h_a}=o(1).
 \label{eq:unsafe-smoothing-mass}
\end{equation}
The decay \(\beta>1/2\) justifies replacing conormal-area
products in this flux as well.  Uniform metric convergence gives
the quadratic-form comparison.  Finally \(H_h(S)<0\) by
\eqref{eq:compact-mean-curvature}; the smooth convergence of the
correction near this disjoint compact surface preserves that
strict inequality.  Each smallness requirement here is imposed
after the fixed solution has been selected.
\end{proof}

\subsection{Enclosures whose boundaries avoid the unsafe set}

\begin{lemma}[Free enclosing boundaries]
\label{lem:free-enclosure}
Choose a sufficiently large fixed scaled radius \(R_*\), then
\(M\) sufficiently large depending on that radius, and perform
the smoothing of Lemma~\ref{lem:corner-smoothing} sufficiently
closely.  There is a perimeter minimizing enclosure \(E\) of
\(O\) and the marked balls about all points of \(\mathcal U_M\),
whose boundary is smooth, free, and disjoint from those balls and
from \(S\).  There is also a largest perimeter minimizing
enclosure \(E_0\) of \(O\) alone.  These sets have all the
boundary, exterior, and minimizing properties in
Proposition~\ref{prop:af-comparison}.
\end{lemma}

\begin{proof}
Extend the smooth inner base geometry across \(S\) in a fixed
collar for the purpose of defining balls.  Mark every closed base
ball
\begin{equation}
 \overline B_{g_{\rm ext}}(P,R_*r_P),\qquad P\in\mathcal U_M,
 \label{eq:unsafe-marked-balls}
\end{equation}
and include their union together with \(O\) as an obstacle.
For fixed parameters the union is compact: the center set is
compact and its positive radius function is continuous.  By
Lemma~\ref{lem:unsafe-blowup}, for each fixed \(R_*\) these
balls and the fixed larger multiples used below lie off the seam
once \(M\) is sufficiently large.  All lie inside the comparator
\(\mathcal S_1\).

Minimize \(h'\)-perimeter among precompact finite-perimeter
sets containing this obstacle modulo null sets.  A minimizer
exists by BV compactness and lower semicontinuity inside a
sufficiently large smooth truncation.  This is also a global
minimizer in the stated class.  To see that the outer truncation
causes no restriction, all sufficiently far AF spheres have
strictly positive outward mean curvature.  The unit normal field
of their foliation has positive divergence.  For almost every
cutting radius, Gauss--Green on the removed part of a competitor
shows that cutting inward saves perimeter if the removed volume
is positive.  The same argument excludes contact with the outer
constraint.  The fixed comparator and the small metric change
give a uniform bound, after enlarging \(C_*\) once,
\begin{equation}
 \Per_{h'}(E)\le C_*.
 \label{eq:unsafe-enclosure-bound}
\end{equation}

We show that its perimeter support cannot meet any marked ball.
Suppose it meets the ball marked at \(P\), at a point \(q\).
Use scaled smooth base coordinates with length unit \(r_P\)
on a fixed multiple of radius \(R_*\).  The base metric and
the comparison metric are equivalent to Euclidean metrics with
constants independent of \(R_*\), by
Lemma~\ref{lem:unsafe-blowup} and the chosen small smoothing.
If this region reaches \(S\), reflect in a base collar for these
comparisons.  Mixed entries of the graph metric vanish on the
face.  No comparison for the arbitrary cap metric is needed:
perimeters of sets containing \(O\) are independent of its
interior extension, and a reflected continuous metric can be
used there in the comparison argument.

The complement of \(E\) has uniform lower volume density at
\(q\) in these scaled coordinates.  Indeed filling a coordinate
ball preserves the obstacle condition.  Minimality and slicing
therefore bound its interior perimeter by a constant times the
complement's trace area on the ball.  If
\(V(a)=|B_a(q)\setminus E|\), Euclidean isoperimetry gives
\begin{equation}
 V(a)^{2/3}\le C V'(a)
 \quad\hbox{for almost every }a,\qquad
 V(a)\ge ca^3.
 \label{eq:unsafe-complement-density}
\end{equation}
The second inequality follows by integrating the first from zero;
the volume is positive at every radius because \(q\) is in
the perimeter support.  Apply this on a ball of radius comparable
to \(R_*\).  On a slightly larger common coordinate ball,
\(E\) itself has volume at least \(cR_*^3\), since it contains
the whole marked base ball modulo null sets.  The uniform
comparison of the scaled base metric with Euclidean metric
justifies this statement even when the marked ball crosses
\(S\).  Relative isoperimetry on that common ball now gives
\begin{equation}
 \Per_{h'}(E)\ge cR_*^2.
 \label{eq:unsafe-obstacle-perimeter}
\end{equation}
Choose \(R_*\) so that this exceeds \(C_*\), and only then
choose \(M\) large enough for all the scaled charts and their
fixed multiples.  This contradicts
\eqref{eq:unsafe-enclosure-bound}.  Compactness of the center set
also covers contacts with the closure of the union of marked balls.
This is where the radius-independent constant in
\eqref{eq:unsafe-quantifiers} is essential.

There is no contact with \(S\) either.  The parallel outward
normal field in a sufficiently thin collar of \(O\) has strictly
negative divergence, since \(H_{h'}(S)<0\).  Adjoin that collar
to a purported minimizer.  Gauss--Green on the added portion
bounds the newly exposed leaf area minus the removed perimeter
by the integral of this negative divergence.  The added volume
is positive for every sufficiently small width at a contact
point; otherwise that point would not be in the perimeter
support.  Slicing allows an almost every width for which all
perimeter formulas hold.  The resulting strict decrease
contradicts minimality.  Adjoining a collar is permitted whether
or not marked balls are present.

The perimeter support is thus separated from both obstacles and
the outer constraint.  It is locally perimeter minimizing in a
smooth three-manifold.  Interior codimension-one minimizing
boundary regularity gives a smooth compact embedded minimal
surface.  Take its open set representative for \(E\).  Filling
any bounded open exterior component would strictly reduce
perimeter, so its open exterior is connected.  Since every
precompact enlargement is an admissible competitor, \(E\) is
outward minimizing.  Its boundary is disjoint from every marked
ball, so the unsafe set lies in its interior.  The metric on
its closed exterior has nonnegative scalar curvature, and its
completeness and unique AF end follow from smoothness of compact
truncations and the AF estimates.

Now minimize with obstacle \(O\) alone.  The same compactness,
outer barriers, and strict inner mean-curvature barrier give
a smooth free minimal boundary and a connected open exterior.
Every component of this open minimizer meets \(O\), because
an open component disjoint from \(O\) can be removed, strictly
reducing perimeter.  Choose a minimizer of maximal volume within
the fixed truncation.  It is a largest minimizer modulo null
sets: the perimeter submodularity inequality implies that the
union and intersection of any two minimizers are minimizers,
and maximal volume forces the union to agree with this one.
Denote its open representative by \(E_0\).  Any proper
precompact enlargement of equal perimeter would also be a
minimizer and would contradict maximal volume.  The outer
barriers allow this reasoning without the truncation.  Thus
\(E_0\) is strictly outward minimizing.  Its boundary avoids
\(S\), so \(O\Subset E_0\).  Finally \(E\) is itself an
admissible enclosure of \(O\), proving
\eqref{eq:unsafe-two-perimeters}.
\end{proof}

No bound on the number or topology of components created by the
marked balls has been used.  In particular the subsequent area
transfer starts from \(E_0\), not from a topological assertion
about \(E\).  The hypotheses for the AF Riemannian Penrose theorem
are satisfied by the closed exterior of \(E\): it is smooth,
complete and one-ended, its scalar curvature is nonnegative and
decays as \(O(|x'|^{-2-2\beta})\), with \(2+2\beta>3\),
its metric decay has exponent \(\beta>1/2\),
and its entire possibly disconnected minimal boundary is outer
minimizing.  The boundary version of \citet[Theorem~1.4]{BrayLee2009}
therefore gives
\begin{equation}
 m_{h'}\ge\sqrt{\frac{\Per_{h'}(E)}{16\pi}}.
 \label{eq:unsafe-af-penrose}
\end{equation}

\subsection{Bad volume and completion of the comparison}

We finish the proof of Proposition~\ref{prop:af-comparison} by
checking the uniform quantitative inputs for
Proposition~\ref{prop:area-transfer}.  The exterior capacity
comparison in Section~\ref{sec:compact-construction} gives
\(x\ge1-\psi_\infty\), so for fixed \(\delta>0\) the set
\(\{x<1-\delta\}\) does not reach beyond a fixed outer sphere.
On the designated region, \eqref{eq:compact-bad-volume} gives
\(\int_{\{x<1-\delta\}}W\le C/L\).  The omitted exterior
collar has \(W=1\), so its contribution is bounded by its base
volume.  The omitted inner collar has base volume \(O(M^{-4})\).
Cauchy--Schwarz and \eqref{eq:compact-height-energy} together
with \eqref{eq:compact-bounds} therefore give
\begin{equation}
 \int_{\{x<1-\delta\}}W\,\dd V_{\rm base}
 \le \frac CL+\Vol_{g_o}(\text{omitted exterior collar})
       +C(L)M^{-1}.
 \label{eq:unsafe-total-bad-volume}
\end{equation}
The coefficient in \(C/L\) is independent of \(L\).  The
other constants may depend on the earlier collar choice and
on \(L\), which is harmless: first make the exterior collar
thin and \(L\) large, then let \(M\) be large.

Choose \(b_*,\delta_0\) small enough in
\eqref{eq:unsafe-mass-limsup} for the prescribed mass tolerance.
For the prescribed \(\delta\) and volume tolerance choose the
collar and \(L\) as just described.  Fix the common comparator
bound, choose \(R_*\) in Lemma~\ref{lem:free-enclosure}, and
then choose \(M\) large enough for that lemma, the scalar
lower bound, the mass costs, and
\eqref{eq:unsafe-total-bad-volume}.  The subsequent small lapse
cutoff and finite strengths are those of
Proposition~\ref{prop:compact-solutions}; all uniform estimates
above hold for those choices.  Apply
Lemma~\ref{lem:corner-smoothing} with mass and metric errors
small enough for the remaining tolerances.  These choices prove
\eqref{eq:unsafe-comparison-output}--\eqref{eq:unsafe-bad-volume-output}.

The caps may be fixed once as topological handlebodies, regardless
of all metric parameters.  The topology entering the area-transfer
argument is consequently the fixed finite topology of
\(X\setminus\operatorname{int}O\).  Positive-time flow levels
outside \(E_0\) are allowed to enter the unsafe region: the
uniform lower scalar bound there is precisely the input supplied
by \eqref{eq:unsafe-comparison-output}.  The mass theorem is
applied to the unmodified smooth exterior of \(E\), whereas
the small initial mean-curvature change used solely for the flow
is covered by Proposition~\ref{prop:area-transfer}.  This
completes the AF comparison construction.
